\documentclass[10pt,a4paper]{amsart}
\usepackage{amsmath,amssymb,amsthm,graphicx}
\usepackage{verbatim}
\newenvironment{reference}{\comment}{\endcomment}
%\newenvironment{reference}{}{}
\newenvironment{slogan}{\comment}{\endcomment}
\newenvironment{history}{\comment}{\endcomment}
% For commutative diagrams we use Xy-pic
\usepackage[all]{xy}

% We use 2cell for 2-commutative diagrams.
\xyoption{2cell}
% Theorem environments.
%
\theoremstyle{plain}
\newtheorem{theorem}[subsection]{Theorem}
\newtheorem{proposition}[subsection]{Proposition}
\newtheorem{lemma}[subsection]{Lemma}

\theoremstyle{definition}
\newtheorem{definition}[subsection]{Definition}
\newtheorem{example}[subsection]{Example}
\newtheorem{exercise}[subsection]{Exercise}
\newtheorem{situation}[subsection]{Situation}

\theoremstyle{remark}
\newtheorem{remark}[subsection]{Remark}
\newtheorem{remarks}[subsection]{Remarks}

\numberwithin{equation}{subsection}

% Macros
%
\def\lim{\mathop{\mathrm{lim}}\nolimits}
\def\colim{\mathop{\mathrm{colim}}\nolimits}
\def\Spec{\mathop{\mathrm{Spec}}}
\def\Hom{\mathop{\mathrm{Hom}}\nolimits}
\def\Ext{\mathop{\mathrm{Ext}}\nolimits}
\def\SheafHom{\mathop{\mathcal{H}\!\mathit{om}}\nolimits}
\def\SheafExt{\mathop{\mathcal{E}\!\mathit{xt}}\nolimits}
\def\Sch{\mathit{Sch}}
\def\Mor{\mathop{\mathrm{Mor}}\nolimits}
\def\Ob{\mathop{\mathrm{Ob}}\nolimits}
\def\Sh{\mathop{\mathit{Sh}}\nolimits}
\def\NL{\mathop{N\!L}\nolimits}
\def\CH{\mathop{\mathrm{CH}}\nolimits}
\def\proetale{{pro\text{-}\acute{e}tale}}
\def\etale{{\acute{e}tale}}
\def\QCoh{\mathit{QCoh}}
\def\Ker{\mathop{\mathrm{Ker}}}
\def\Im{\mathop{\mathrm{Im}}}
\def\Coker{\mathop{\mathrm{Coker}}}
\def\Coim{\mathop{\mathrm{Coim}}}

% Boxtimes
%
\DeclareMathSymbol{\boxtimes}{\mathbin}{AMSa}{"02}

%
% Macros for moduli stacks/spaces
%
\def\QCohstack{\mathcal{QC}\!\mathit{oh}}
\def\Cohstack{\mathcal{C}\!\mathit{oh}}
\def\Spacesstack{\mathcal{S}\!\mathit{paces}}
\def\Quotfunctor{\mathrm{Quot}}
\def\Hilbfunctor{\mathrm{Hilb}}
\def\Curvesstack{\mathcal{C}\!\mathit{urves}}
\def\Polarizedstack{\mathcal{P}\!\mathit{olarized}}
\def\Complexesstack{\mathcal{C}\!\mathit{omplexes}}
% \Pic is the operator that assigns to X its picard group, usage \Pic(X)
% \Picardstack_{X/B} denotes the Picard stack of X over B
% \Picardfunctor_{X/B} denotes the Picard functor of X over B
\def\Pic{\mathop{\mathrm{Pic}}\nolimits}
\def\Picardstack{\mathcal{P}\!\mathit{ic}}
\def\Picardfunctor{\mathrm{Pic}}
\def\Deformationcategory{\mathcal{D}\!\mathit{ef}}
\usepackage{cleveref}
\crefname{originalsection}{original section}{original sections}
\Crefname{originalsection}{Original section}{Original sections}
\makeatletter
\newlabel{section-conventions}{{section-conventions}{source}{}{sourcepointer-1}{}}
\newlabel{section-conventions@cref}{{[originalsection][0][]section-conventions}{[1][0][]source}{}{}{}}
\makeatother
\crefname{originalreference}{original reference}{original references}
\Crefname{originalreference}{Original reference}{Original references}
\makeatletter
\newlabel{sites-lemma-mono-epi-sheaves}{{sites-lemma-mono-epi-sheaves}{source}{}{sourcepointer-2}{}}
\newlabel{sites-lemma-mono-epi-sheaves@cref}{{[originalreference][0][]sites-lemma-mono-epi-sheaves}{[1][0][]source}{}{}{}}
\makeatother
\crefname{originalreference}{original reference}{original references}
\Crefname{originalreference}{Original reference}{Original references}
\makeatletter
\newlabel{stacks-lemma-when-gerbe}{{stacks-lemma-when-gerbe}{source}{}{sourcepointer-3}{}}
\newlabel{stacks-lemma-when-gerbe@cref}{{[originalreference][0][]stacks-lemma-when-gerbe}{[1][0][]source}{}{}{}}
\makeatother
\title{A gerbe over an algebraic space is the sheafification of the presheaf of isomorphism classes of its objects}
\author{The Stacks Project Authors}
\date{stacks-project:stacks-morphisms.tex:local-snapshot-20261012; GNU FDL 1.2 or later}
\begin{document}
\maketitle
\noindent\textit{Editorial note.} All delivered mathematical text is the author's own, copied from the pinned Stacks Project chapter source sources/stacks-algebraic-stacks-morphisms/pinned-original-source.tex; no statement or proof step was rewritten, reordered or removed.\par\medskip
\noindent\textit{Editorial note.} Preamble provenance: the Stacks Project keeps one single shared preamble for every chapter; this card takes it directly from the staged sources/stacks-algebraic-stacks-morphisms/preamble.tex as three verbatim spans (lines 8-12, 14-18 and 154-217), so no derived file was written and nothing was added or redefined. The wrapper class and the base packages are supplied by the assembler.\par\medskip
\noindent\textit{Editorial note.} Bibliography provenance: neither the counted statement nor its proof contains a citation (this lane rescans the whole span before writing the card), so the card carries no bibliography and no entry was reconstructed.\par\medskip
\noindent\textit{Editorial note.} Selection provenance: this card was opened under the file-based control contract (CONTROL.md), with the mechanical pre-card gate run on its statement and proof interval before the card was written.\par\medskip
\noindent\textit{Allowed standard prerequisites (editorial).}\begin{itemize}
\item The two clauses of the gerbe condition: fppf local existence of objects and fppf local existence of isomorphisms. Invoked twice inside the counted proof (source lines 6110-6113 and 6119-6122). Supplies the local objects needed for surjectivity and the local isomorphisms needed for injectivity.
\item The criterion that a morphism of sheaves which is injective and surjective is an isomorphism. Invoked inside the counted proof (source lines 6114-6116) as a cross-chapter pointer. Upgrades the two computations to the asserted isomorphism of sheaves.
\end{itemize}\medskip
\ifdefined\hypertarget\hypertarget{sourcepointer-1}{}\fi
\noindent\textit{Source reference.} Original section section-conventions: sources/stacks-algebraic-stacks-morphisms/pinned-original-source.tex, source lines 257--257. Navigation pointer only: the counted text refers to the Stacks item section-conventions, which is an existing labeled result of this same chapter and is not part of this card. The pointer carries the target's own label as its number and its definition line as its span; its content is not reproduced here.\par\medskip
\ifdefined\hypertarget\hypertarget{sourcepointer-2}{}\fi
\noindent\textit{Source reference.} Original reference sites-lemma-mono-epi-sheaves: sources/stacks-algebraic-stacks-morphisms/pinned-original-source.tex, source lines 6114--6114. Navigation pointer only: the counted text refers to the Stacks item sites-lemma-mono-epi-sheaves, which belongs to another chapter of the project and is not part of this card. The pointer carries the target's own label as its number because this chapter's staged snapshot does not contain the other chapter's tag table; its content is not reproduced here.\par\medskip
\ifdefined\hypertarget\hypertarget{sourcepointer-3}{}\fi
\noindent\textit{Source reference.} Original reference stacks-lemma-when-gerbe: sources/stacks-algebraic-stacks-morphisms/pinned-original-source.tex, source lines 6111--6111. Navigation pointer only: the counted text refers to the Stacks item stacks-lemma-when-gerbe, which belongs to another chapter of the project and is not part of this card. The pointer carries the target's own label as its number because this chapter's staged snapshot does not contain the other chapter's tag table; its content is not reproduced here.\par\medskip
\begin{lemma}
\label{lemma-gerbe-over-iso-classes}
Let $\mathcal{X}$ be an algebraic stack. If $\mathcal{X}$ is a gerbe, then
the sheafification of the presheaf
$$
(\Sch/S)_{fppf}^{opp} \to \textit{Sets}, \quad
U \mapsto \Ob(\mathcal{X}_U)/\!\!\cong
$$
is an algebraic space and $\mathcal{X}$ is a gerbe over it.
\end{lemma}

\begin{proof}
(In this proof the abuse of language introduced in
Section \ref{section-conventions}
really pays off.)
Choose a morphism $\pi : \mathcal{X} \to X$ where $X$ is an
algebraic space which turns $\mathcal{X}$ into a gerbe over $X$.
It suffices to prove that $X$ is the sheafification of the presheaf
$\mathcal{F}$ displayed in the lemma.
It is clear that there is a map $c : \mathcal{F} \to X$.
We will use
Stacks, Lemma \ref{stacks-lemma-when-gerbe}
properties (2)(a) and (2)(b) to see that the map $c^\# : \mathcal{F}^\# \to X$
is surjective and injective, hence an isomorphism, see
Sites, Lemma \ref{sites-lemma-mono-epi-sheaves}.
Surjective: Let $T$ be a scheme and let $f : T \to X$. By property (2)(a)
there exists an fppf covering $\{h_i : T_i \to T\}$ and morphisms
$x_i : T_i \to \mathcal{X}$ such that $f \circ h_i$ corresponds to
$\pi \circ x_i$. Hence we see that $f|_{T_i}$ is in the image of $c$.
Injective: Let $T$ be a scheme and let $x, y : T \to \mathcal{X}$
be morphisms such that $c \circ x = c \circ y$.
By (2)(b) we can find a covering $\{T_i \to T\}$ and morphisms
$x|_{T_i} \to y|_{T_i}$ in the fibre category $\mathcal{X}_{T_i}$.
Hence the restrictions $x|_{T_i}, y|_{T_i}$ are equal in
$\mathcal{F}(T_i)$. This proves that $x, y$ give the same section
of $\mathcal{F}^\#$ over $T$ as desired.
\end{proof}

\end{document}
